% Part: incompleteness
% Chapter: representability-in-q
% Section: c

\documentclass[../../../include/open-logic-section]{subfiles}

\begin{document}

\olfileid{inc}{req}{cee}
\olsection{फलनांचा संच~$C$}

$C$ हा पुढील फलने समाविष्ट करणारा सर्वात लहान फलनसंच असू द्या:
\begin{enumerate}
\item $0$,
\item उत्तरवर्ती फलन,
\item बेरीज,
\item गुणाकार,
\item प्रक्षेपण फलने, आणि
\item समतेचे जातिबोधक फल, $\Char{=}$;
\end{enumerate}
आणि जो पुढील क्रियांखाली बंद आहे:
\begin{enumerate}
\item संयोजन, आणि
\item नियमित फलनांवर लावलेला अपरिबद्ध शोध.
\end{enumerate}
ही शेवटची मर्यादा म्हणजे $\mu$ क्रिया फक्त तिचा परिणाम
सर्वत्र परिभाषित असेल तेव्हाच वापरता येते, एवढेच. याची
\emph{सामान्य पुनरावर्ती} फलनांच्या व्याख्येशी तुलना करा:
येथे बेरीज, गुणाकार आणि $\Char{=}$ समाविष्ट केले आहेत,
पण आदिम पुनरावर्तन वगळले आहे.

$C$ मधील प्रत्येक फलन पुनरावर्ती आहे, हे स्पष्ट आहे; कारण
बेरीज, गुणाकार आणि $\Char{=}$ ही फलने पुनरावर्ती आहेत.
याचा उलट दावा खरा आहे, हेही दाखवू. म्हणजे $C$ मधील
इतर क्रिया वापरून आदिम पुनरावर्तन करता येते.

\end{document}
